\FloatBarrier
{\color{blue}\section{Conclusions and Outlook}
\label{sec:Conclusions}}

\GPT{Suggested points: summarize the distinction between trajectory-resolved critical chiral boundaries and the short-range outcome average; separate the evidence for criticality and chirality from the open identification with an equilibrium ground state.}

\claudeadd{We have shown that repeated finite-range occupation measurements with purely local feedback prepare chiral critical boundaries without postselection. In a slab geometry, the circuit prepares a Chern bulk, and each measurement record yields a pure state whose interfaces carry algebraic correlations, logarithmic entanglement and charge fluctuations with coefficients close to those of a free chiral fermion, half a unit of entropy coefficient per wall, and opposite modular handedness on the two walls. The same walls host the slowest modes of the purification dynamics. All of these signatures are properties of the record-resolved ensemble. Discarding the record leaves a mixed state with short-range correlations and a gapped averaged channel, consistent with the obstructions to dissipative preparation of Chern insulators. Whether the conditioned wall states are equivalent to the ground state of an equilibrium chiral Luttinger liquid remains open.}

\claudeadd{Because the trajectory-level criticality is invisible in the averaged state, observing it experimentally faces the postselection problem of monitored dynamics. Here, the Gaussian structure helps: each record's state can be computed classically from the measured outcomes, which enables postselection-free estimators that combine single-copy measurements with classical simulation~\cite{GarrattAltman2024,McGinley2024}.}
\claudecom{On the outlook items: (i) observability, as in the paragraph above; (ii) depth: whether a layered schedule of commuting measurements, started from a product state, reaches the critical walls in a number of cycles independent of $N_y$, given that the Lyapunov gap suggests relaxation times growing with $N_y$; (iii) class D and vortices with Majorana zero modes, as GPT suggests, which would need a different stabilizer construction; (iv) robustness to imperfect feedback or interactions. Pick two; four is too many for an outlook.}

\GPT{Possible extensions include other symmetry classes and programmable point defects, such as vortices hosting Majorana zero modes in two-dimensional superconducting systems. This is an outlook direction, not a result established here.}


